% META: {"id": "TSPE-SUITLIM-EX-046", "chapitre": "TSPE-SUITES-LIMITES", "type_objet": "exercice", "capacites": ["TSPE-SUITLIM-C7"], "capacites_codes": ["C7"], "methodes": ["M5"], "parcours": 2, "competences": ["raisonner", "calculer"], "duree_min": 15, "mode_creation": "ex_nihilo", "sources_inspiration": [], "fichier_tex": "chapitres/TSPE-SUITES-LIMITES/exercices/TSPE-SUITLIM-EX-046.tex", "corrige_tex": "chapitres/TSPE-SUITES-LIMITES/corriges/TSPE-SUITLIM-CO-046.tex", "status": "generated"}
% BEGIN-VERIFY
% from sympy import *
% x = symbols('x')
% f = (2*x + 1)*exp(-x)
% assert limit(f, x, oo) == 0
% assert limit(f, x, -oo) == -oo
% fp = diff(f, x)
% # f'(x) = 2*exp(-x) - (2x+1)*exp(-x) = exp(-x)*(1-2x)
% assert simplify(fp - exp(-x)*(1-2*x)) == 0
% END-VERIFY
\begin{exercice}{TSPE-SUITLIM-EX-046}{2}{15}
Soit $f(x) = (2x + 1)\,\mathrm{e}^{-x}$.
\begin{enumerate}
  \item Determiner $\lim_{x \to +\infty} f(x)$ et $\lim_{x \to -\infty} f(x)$.
  \item Calculer $f'(x)$ et etudier les variations de $f$ sur $\mathbb{R}$.
  \item La droite $y = 0$ est-elle asymptote horizontale ? Justifier.
\end{enumerate}
\end{exercice}
